\documentclass{article}
\usepackage{lmodern}
\usepackage{amsmath}
\usepackage{listings}
\usepackage{fullpage}
\usepackage{parskip}
\usepackage{xcolor}
\lstset{
	basicstyle=\ttfamily,
	columns=fullflexible,
	frame=single,
	breaklines=true,
	postbreak=\mbox{\textcolor{red}{$\hookrightarrow$}\space},
}
\begin{document}
\title{Experiment `p\_6\_9\_arrays\_equal\_array' Results}
\author{\today}
\date{}
\maketitle
\textbf{Experiment outcome: }SUCCESS
\\\textbf{Bad responses: }0
\\\textbf{Responses containing}~\texttt{assume}~\textbf{: }0
\\\textbf{Responses containing}~\texttt{expect}~\textbf{: }0
\\\textbf{Resolution attempts: }1
\\\textbf{Hard fails (resolution): }0
\\\textbf{Soft fails (resolution): }0
\\\textbf{Verification attempts: }1
\section*{Problem Specification}
\textbf{Problem name: }p\_6\_9\_arrays\_equal\_array
\\\textbf{Natural language statement: }Write a method that checks whether two arrays have the same elements in the same order.
\\\textbf{Method signature: }p\_6\_9\_arrays\_equal\_array(first: array<int>, second: array<int>) returns (are\_equal: bool)
\subsection*{Ensures}
\begin{itemize}
\item \texttt{are\_equal == (first.Length == second.Length \&\& forall i :: 0 <= i < first.Length ==> first[i] == second[i])}
\end{itemize}
\clearpage
\section*{GenAI interactions}
Below you will find all interactions between the `user' (program) and the `assistant' (OpenAI).
\subsection*{Program $\rightarrow$ GenAI}
\begin{lstlisting}
You are given the following task to perform in Dafny:

Write a method that checks whether two arrays have the same elements in the same order.

The signature should be:

p_6_9_arrays_equal_array(first: array<int>, second: array<int>) returns (are_equal: bool)

The method should respect the following contract:

ensures are_equal == (first.Length == second.Length && forall i :: 0 <= i < first.Length ==> first[i] == second[i])

Produce and show only the Dafny body of this method, including the curly braces that surround it. Do not show the signature nor contract. You must not use 'assume' nor 'expect' anywhere in your code.
\end{lstlisting}
\subsection*{GenAI $\rightarrow$ Program}
\textbf{System fingerprint: }null
\\\textbf{ID: }chatcmpl-Dn1jD2AcgOc4scarbytaidXmSYxi4
\\\textbf{Model: }gpt-5.5-2026-04-23
\\\textbf{Created at: }1780577087
\\\textbf{Finish reason: }stop
\\\textbf{Usage: }320 tokens in, and 3087 tokens out
\begin{lstlisting}
{
  if first.Length != second.Length {
    are_equal := false;
    assert !(first.Length == second.Length);
    assert !(first.Length == second.Length && (forall i :: 0 <= i < first.Length ==> first[i] == second[i]));
  } else {
    var i := 0;
    are_equal := true;
    while i < first.Length
      invariant first.Length == second.Length
      invariant 0 <= i <= first.Length
      invariant are_equal ==> (forall k :: 0 <= k < i ==> first[k] == second[k])
      invariant !are_equal ==> (exists k :: 0 <= k < i && first[k] != second[k])
      decreases first.Length - i
    {
      if first[i] != second[i] {
        are_equal := false;
        assert 0 <= i < i + 1 && first[i] != second[i];
        assert exists k :: 0 <= k < i + 1 && first[k] != second[k];
      } else {
        assert first[i] == second[i];
        if are_equal {
          assert forall k :: 0 <= k < i ==> first[k] == second[k];
          forall k | 0 <= k < i + 1
            ensures first[k] == second[k]
          {
            if k < i {
              assert first[k] == second[k];
            } else {
              assert k == i;
            }
          }
        } else {
          ghost var k :| 0 <= k < i && first[k] != second[k];
          assert 0 <= k < i + 1 && first[k] != second[k];
          assert exists j :: 0 <= j < i + 1 && first[j] != second[j];
        }
      }
      i := i + 1;
    }

    assert i == first.Length;
    if are_equal {
      assert forall k :: 0 <= k < first.Length ==> first[k] == second[k];
    } else {
      ghost var k :| 0 <= k < first.Length && first[k] != second[k];
      if forall j :: 0 <= j < first.Length ==> first[j] == second[j] {
        assert first[k] == second[k];
        assert false;
      }
      assert !(forall j :: 0 <= j < first.Length ==> first[j] == second[j]);
    }
  }
}
\end{lstlisting}
\section*{Final Program}
\begin{lstlisting}
method p_6_9_arrays_equal_array(first: array<int>, second: array<int>) returns (are_equal: bool)
	ensures are_equal == (first.Length == second.Length && forall i :: 0 <= i < first.Length ==> first[i] == second[i])
{
  if first.Length != second.Length {
    are_equal := false;
    assert !(first.Length == second.Length);
    assert !(first.Length == second.Length && (forall i :: 0 <= i < first.Length ==> first[i] == second[i]));
  } else {
    var i := 0;
    are_equal := true;
    while i < first.Length
      invariant first.Length == second.Length
      invariant 0 <= i <= first.Length
      invariant are_equal ==> (forall k :: 0 <= k < i ==> first[k] == second[k])
      invariant !are_equal ==> (exists k :: 0 <= k < i && first[k] != second[k])
      decreases first.Length - i
    {
      if first[i] != second[i] {
        are_equal := false;
        assert 0 <= i < i + 1 && first[i] != second[i];
        assert exists k :: 0 <= k < i + 1 && first[k] != second[k];
      } else {
        assert first[i] == second[i];
        if are_equal {
          assert forall k :: 0 <= k < i ==> first[k] == second[k];
          forall k | 0 <= k < i + 1
            ensures first[k] == second[k]
          {
            if k < i {
              assert first[k] == second[k];
            } else {
              assert k == i;
            }
          }
        } else {
          ghost var k :| 0 <= k < i && first[k] != second[k];
          assert 0 <= k < i + 1 && first[k] != second[k];
          assert exists j :: 0 <= j < i + 1 && first[j] != second[j];
        }
      }
      i := i + 1;
    }

    assert i == first.Length;
    if are_equal {
      assert forall k :: 0 <= k < first.Length ==> first[k] == second[k];
    } else {
      ghost var k :| 0 <= k < first.Length && first[k] != second[k];
      if forall j :: 0 <= j < first.Length ==> first[j] == second[j] {
        assert first[k] == second[k];
        assert false;
      }
      assert !(forall j :: 0 <= j < first.Length ==> first[j] == second[j]);
    }
  }
}
\end{lstlisting}
\section*{Total Token Usage}
\textbf{Input tokens: }320
\\\textbf{Output tokens: }3087
\\\textbf{Reasoning tokens: }2486
\\\textbf{Sum of `total tokens': }3407
\section*{Experiment Timings}
\textbf{Overall Experiment} started at 1780577084225, ended at 1780577168851, lasting 84626ms (84.63 seconds)
\\\textbf{Iteration \#1} started at 1780577084229, ended at 1780577168851, lasting 84622ms (84.62 seconds)
\end{document}
